\documentclass{article}
\usepackage[T1]{fontenc}
\usepackage{lmodern}
\usepackage{graphicx}
\usepackage{amsmath,amssymb}
\usepackage[a4paper,margin=2cm]{geometry}
\usepackage[absolute,overlay]{textpos}
\setlength{\TPHorizModule}{1mm}
\setlength{\TPVertModule}{1mm}
\usepackage[indonesian]{babel}
\usepackage{hyperref}
\setlength{\emergencystretch}{2em}
% unit: Halaman 76

\begin{document}

% === Halaman 76 (python/native) ===

Inisialisasi:

• Masukkan 80-bit initialization vector (IV) ke dalam A

• Masukkan 80-bit kunci ke dalam B

• Set c109 , c110 , c111 =1, sedangkan bit lainnya 0

Pemanasan (warming-up):

• Detak (clock) NLFSR  4 x 288 = 1152 kali tanpa menghasilkan luaran (luaran diabaikan)

• Tujuannya adalah melakukan pencampuran kunci/IV sehingga pengaruh kunci dan IV tersebar ke seluruh register

Encryption:

• XOR-kan ketiga buah luaran NLFSR untuk membangkitkan kunci-alir si 

• XOR-kan si dengan bit plainteks pi

76

% deskripsi: gambar native xref 304
\noindent\includegraphics[width=0.85\textwidth]{../../images/python/page_076_img_01.png}

\end{document}
